\documentclass[12pt]{article}

% ------------------------------------------------------------
% LuaLaTeX Core Packages
% ------------------------------------------------------------

\usepackage{fontspec}
\setmainfont{Libertinus Serif}

\usepackage{microtype}
\usepackage{geometry}
\geometry{margin=1in}

\usepackage{setspace}
\onehalfspacing

% ------------------------------------------------------------
% Mathematics
% ------------------------------------------------------------

\usepackage{amsmath}
\usepackage{amssymb}
\usepackage{mathtools}

% ------------------------------------------------------------
% Theorem Environments
% ------------------------------------------------------------

\usepackage{amsthm}

\newtheorem{definition}{Definition}
\newtheorem{principle}{Principle}
\newtheorem{proposition}{Proposition}
\newtheorem{lemma}{Lemma}
\newtheorem{theorem}{Theorem}

% ------------------------------------------------------------
% Diagrams
% ------------------------------------------------------------

\usepackage{tikz}
\usepackage{tikz-cd}

% ------------------------------------------------------------
% Hyperlinks
% ------------------------------------------------------------

\usepackage{hyperref}

\hypersetup{
colorlinks=true,
linkcolor=black,
citecolor=black,
urlcolor=blue
}

% ------------------------------------------------------------
% Title Information
% ------------------------------------------------------------

\title{Execution, History, and Reduction:\\
A Minimal Operational Kernel for Computation}

\author{Flyxion}

\date{\today}

\begin{document}

\maketitle

\begin{abstract}

Most computational formalisms describe systems in terms of
state transformations. Programs update states, machines transition
between configurations, and system behavior is modeled as a sequence
of state changes.

This paper proposes an alternative foundation in which the primary
objects of computation are not states but histories. A history records
the irreversible accumulation of events and therefore encodes the
constraints that those events impose on future execution.

Within this perspective, computation is generated by a minimal
operational kernel consisting of three primitive operations:
extension, merge, and reduction. Extension appends new events to an
existing history, producing longer executions. Merge reconciles
compatible histories by identifying shared prefixes and combining
their subsequent developments. Reduction removes distinctions between
histories by collapsing them into coarser representations.

These operations induce an ordered structure over histories in which
more detailed histories refine less detailed ones. The resulting space
of executions forms an order-enriched monoidal category whose
composition laws capture the fundamental algebra of computation.

This history-based perspective clarifies structural properties that
are difficult to express in traditional state-transition models.
In particular, the existence of lawful intermediate execution steps
(divisibility) becomes a property of history factorization rather
than of state evolution. Dual operational representations acting on
states and observables may impose distinct admissibility constraints
on such factorizations, revealing an intrinsic directional structure
in the algebra of execution.

The minimal operational kernel developed here unifies patterns that
appear across distributed systems, event-sourced architectures,
version-control systems, and dynamical models of computation.
Viewed through this lens, computation emerges as the geometry of
irreversible histories evolving within an ordered space of
possibilities.

\end{abstract}

\tableofcontents

\newpage

% ------------------------------------------------------------
% Introduction
% ------------------------------------------------------------

\section*{Introduction}

The dominant traditions of computer science describe programs
in terms of transformations between states. From early automata
theory to modern programming-language semantics, computational
systems are typically modeled as structures that occupy
well-defined configurations and evolve through transitions
between them. These models have proven extraordinarily
successful for reasoning about correctness, complexity,
and program behavior.

At the same time, the systems that modern computing
actually relies upon increasingly operate in ways that
strain the assumptions of state-centered models.
Large-scale distributed services, append-only logs,
version-control systems, and event-sourced architectures
do not treat state as the primary object of computation.
Instead they record sequences of events whose cumulative
effects determine the observable condition of the system.

In such systems, the present configuration is best understood
not as a self-contained object but as a compressed summary
of the history that produced it. The ordering of events,
the causal relations among them, and the constraints
propagated through their interactions carry information
that cannot be recovered from a snapshot alone.

This observation suggests a different conceptual starting
point for the theory of computation. Rather than treating
states as primary and histories as auxiliary traces,
one may reverse the relationship and treat histories
as the fundamental structure from which states are derived.
On this view, computation can be understood as the algebra
of history composition, generated by extension, merge, and
reduction over an ordered space of executions.

The aim of this essay is to develop a minimal operational
framework built upon this reversal. In this framework,
computation is interpreted as the irreversible construction
of event histories under constraint. System states emerge
as reductions of those histories, and abstraction itself
is understood as a controlled compression of execution.

The following sections develop this perspective step by step.
We begin by examining the limitations of state-centered
ontologies of computation and then introduce an alternative
kernel in which event histories serve as the primitive
structure of computational systems.

% ------------------------------------------------------------
% 1. The Failure of State Ontology
% ------------------------------------------------------------

\section{The Failure of State Ontology}

Most formal models of computation treat programs as transformations
between states. A system is assumed to occupy a well-defined state,
and computation is represented as a sequence of transitions between
such states.

This perspective is historically rooted in set-theoretic semantics.
Programs are interpreted as relations over sets of possible states,
and correctness is defined in terms of membership within these sets.

While this approach has produced powerful theoretical tools,
it imposes an ontological assumption that does not accurately
reflect the behavior of real computational systems.

\subsection{The Snapshot Assumption}

State-based models assume that the complete condition of a system
can be captured by a snapshot taken at a particular moment in time.

In practice, however, many systems cannot be fully described
by instantaneous snapshots. Distributed databases, version-control
systems, and streaming computation engines all rely on persistent
event histories rather than static states.

A snapshot may summarize a system, but it does not explain
how that system came to be.

\subsection{Loss of Causal Structure}

When computation is modeled purely as state transformation,
the causal structure of execution is obscured.

Two states may appear identical while having arisen through
entirely different sequences of events. These distinct histories
can have radically different implications for correctness,
safety, and reproducibility.

For example, the state of a distributed ledger cannot be validated
without examining the history of transactions that produced it.
Similarly, the meaning of a Git repository is not contained in the
current tree but in the commit graph that generated it.

In such systems, history is not auxiliary metadata.
It is the primary structure of the system.

\subsection{Irreversibility of Execution}

Execution is fundamentally irreversible. Once an event occurs,
it cannot be undone without introducing additional compensating
events.

This irreversibility means that the evolution of a system
is more accurately described as the growth of a history
rather than the traversal of a state space.

The state-space model attempts to compress this process into
a sequence of snapshots, but this compression discards
information about causal ordering and construction.

\subsection{Toward an Event-Historical Ontology}

The limitations of state ontology suggest a different
foundational assumption:

\begin{principle}[Event-Historical Primacy]
The primary structure of computation is the irreversible
history of events that construct a system.
\end{principle}

Within this perspective, a state is not a primitive object
but a derived summary produced by replaying a history.

In other words, state becomes a projection of history.

This shift from state to history has significant consequences.
It changes how we understand determinism, composition,
abstraction, and correctness.

The remainder of this essay develops a minimal operational
kernel built upon this premise.

% ------------------------------------------------------------
% 2. Event Histories as the Primitive Structure
% ------------------------------------------------------------

\section{Event Histories as the Primitive Structure}

If computation is fundamentally the construction of irreversible
histories, then the primitive object of the theory must be
the history itself rather than the state it produces.

\subsection{Events}

We begin by defining the atomic unit of execution.

\begin{definition}[Event]
An \emph{event} is a record of an irreversible transition in a system.
Formally, an event is an element
\[
e \in \mathcal{E}
\]
where $\mathcal{E}$ denotes the set of all possible events
in a given computational universe.
\end{definition}

Events represent observable transitions such as message delivery,
state mutation, computation steps, or external inputs.

The internal structure of events is not fixed by the theory;
only their ordering relations are essential.

\subsection{Histories}

A computational system evolves through the accumulation of events.

\begin{definition}[Event History]
An \emph{event history} $H$ is a finite or countably infinite
sequence of events
\[
H = (e_1, e_2, e_3, \dots)
\]
such that each event represents a valid extension
of the history that precedes it.
\end{definition}

The essential property of a history is that it grows monotonically:
events may be appended but not removed.

\begin{principle}[Irreversibility]
If $H$ is a history and $e$ is an event appended to $H$,
then the resulting history
\[
H' = H \cdot e
\]
strictly extends $H$.
\end{principle}

This extension operator defines the fundamental notion of execution.

\subsection{Execution as History Extension}

Execution is therefore defined not as state transition
but as the extension of histories.

\begin{definition}[Execution]
An execution step is the operation
\[
H \rightarrow H \cdot e
\]
which appends an event $e$ to the history $H$.
\end{definition}

Repeated application of this operation constructs
the entire observable evolution of a system.

\subsection{Partial Ordering of Events}

In many systems events are not totally ordered.
Independent operations may occur concurrently.

To capture this, histories can be generalized from sequences
to partially ordered sets.

\begin{definition}[Event Poset]
An event history may be represented as a partially ordered set
\[
(H, \leq)
\]
where $e_i \leq e_j$ indicates that event $e_i$
causally precedes event $e_j$.
\end{definition}

This structure captures concurrency while preserving causal order.

\subsection{Derived State}

State now appears as a derived object.

\begin{definition}[State Projection]
A \emph{state} is a deterministic function
\[
\sigma : H \rightarrow S
\]
that maps an event history $H$ to a representation
in some state space $S$.
\end{definition}

Different projections may summarize different aspects
of the same history.

For example, a database snapshot, a filesystem tree,
or the current contents of memory may all be produced
by applying distinct projections to the same underlying history.

Thus state is not fundamental.
It is a compressed view of execution history.

% ------------------------------------------------------------
% 3. The Minimal Operational Kernel
% ------------------------------------------------------------

\section{The Minimal Operational Kernel}

Having established event histories as the primitive
structure of computation, we may now identify the
minimal operations required to construct and transform
those histories. The objective of this section is to
characterize a small set of operations sufficient to
generate the class of computational systems described
throughout this work.

Within the event-historical framework, execution may
be understood as the progressive construction and
manipulation of histories through three fundamental
transformations. The first transformation extends a
history by the addition of a new event. This operation,
referred to as \emph{extension}, corresponds to the
irreversible execution of a transition whose effects
become incorporated into the growing causal structure
of the system.

The second transformation combines compatible histories
into a common continuation. This operation, referred to
as \emph{merge}, reconciles divergent execution paths
by constructing a history that preserves the causal
ordering of events from both inputs. When defined over
compatible histories, merge corresponds to a join-like
operation in the partially ordered space of histories.

The third transformation derives simplified
representations from existing histories. This operation,
referred to as \emph{reduction}, maps histories into
compressed descriptions that preserve selected
properties while discarding others. Reductions give
rise to states, abstractions, and summaries that allow
systems to reason about their execution without
retaining the full causal trace.

Together, extension, merge, and reduction constitute a
minimal operational kernel for history-based computation.
Extension generates histories through execution, merge
reconciles concurrent or divergent histories, and
reduction produces derived representations suitable
for analysis and control. The interaction of these
operations is sufficient to generate the space of
executable systems considered in this framework.

\subsection{History Extension}

Execution is the irreversible extension of an event history.

\begin{definition}[History Extension]
Let $H$ be a history and $e$ an event. The extension operator
\[
\mathrm{extend}(H,e)
\]
produces a new history
\[
H' = H \cdot e.
\]
\end{definition}

This operation represents the fundamental step of computation.
Every observable transition in a system corresponds to the
extension of its event history.

Because extension is irreversible, histories grow monotonically.

\begin{principle}[Monotonic Growth]
If $H_1$ precedes $H_2$ in execution, then $H_1$ is a prefix
of $H_2$.
\end{principle}

This monotonicity guarantees that execution histories
form a directed structure rather than a reversible trajectory.

\subsection{History Merge}

Independent histories may be combined when their events
are mutually consistent.

\begin{definition}[History Merge]
Let $H_1$ and $H_2$ be histories that share a common prefix.
A merge operation produces a history $H_3$ that contains
all events from both histories while preserving causal order.
\end{definition}

Formally, if
\[
H_1 = H_0 \cdot A,
\qquad
H_2 = H_0 \cdot B,
\]
then the merged history is a causal extension containing
the events of both $A$ and $B$.

Merge therefore constructs a joint continuation of the
two histories.

This operation is central to distributed computation,
version-control systems, and log-based architectures.

\subsection{Compatibility}

Not all histories can be merged.

\begin{definition}[Compatibility]
Two histories are compatible if their events can be combined
without violating causal constraints or system invariants.
\end{definition}

Compatibility acts as the primary constraint on composition.

If histories are incompatible, they cannot be merged directly;
instead compensating events must be introduced.

\subsection{History Reduction}

The final primitive operation is reduction.

Reduction compresses histories by discarding distinctions
that are irrelevant for a given purpose.

\begin{definition}[History Reduction]
A reduction is a mapping
\[
\rho : H \rightarrow H'
\]
that preserves selected structural properties of a history
while discarding others.
\end{definition}

Reductions include a wide range of familiar operations that
derive compact representations from underlying execution
histories. Examples include the generation of snapshots that
summarize the configuration of a system at a particular point
in its evolution, the construction of state representations
that aggregate the effects of prior events, and the formation
of higher-level abstractions that retain only those aspects
of the history relevant for reasoning or control. Aggregation
operations used in monitoring systems, database views, and
analytical pipelines similarly produce condensed summaries of
the underlying execution trace.

Despite their practical importance, such operations do not
introduce new structural information into the system. Instead,
they operate by discarding portions of the causal structure
encoded in the history. Formally, a reduction operator
\[
\rho : \mathcal{H} \rightarrow S
\]
maps histories to derived representations by compressing
their internal structure while preserving selected
observables.

The key observation is therefore that reductions do not
generate new computational structure. They remove information
from existing histories by collapsing multiple distinct
histories into a single representation. In this sense,
reductions act as information-losing morphisms that project
the detailed structure of execution histories into spaces
of simplified descriptions. The structural complexity of the
system remains entirely contained within the histories
themselves, while reductions merely provide compressed views
of that structure.

\subsection{Abstraction as Reduction}

This leads to a reinterpretation of abstraction.

\begin{principle}[Abstraction as Reduction]
An abstraction is a reduction of an event history that preserves
only the distinctions relevant to a particular task.
\end{principle}

Under this view, abstraction is not the creation of higher-level
objects but the controlled loss of information.

\subsection{Kernel Sufficiency}

These three operations—extension, merge, and reduction—are
sufficient to construct the computational structures
observed in real systems.

Execution produces histories.

Composition merges histories.

Abstraction reduces histories.

From these operations arise deterministic replay,
distributed consensus mechanisms,
versioned data structures,
and reproducible computation.

The remainder of the essay explores the consequences
of this minimal kernel.

% ------------------------------------------------------------
% 4. Determinism and Replayability
% ------------------------------------------------------------

\section{Determinism and Replayability}

Once computation is defined in terms of event histories,
determinism can be reformulated in a more precise way.

Traditional models treat determinism as a property of
state-transition functions. A program is deterministic
if the same state always produces the same next state.

Within the event-historical framework, determinism
is instead a property of replay.

\subsection{Deterministic Replay}

\begin{definition}[Replay]
Given a history
\[
H = (e_1, e_2, \dots, e_n),
\]
replay is the process of re-executing the sequence
of events in order to reconstruct the system state.
\end{definition}

Replay reconstructs a system by applying each event
in the order specified by the history.

\begin{definition}[Replay Determinism]
A system is \emph{replay-deterministic} if identical
histories always produce identical derived states.
\end{definition}

In other words, determinism arises when the projection
\[
\sigma : H \rightarrow S
\]
is stable under replay.

\subsection{State as a Derived Cache}

Within this perspective, state is best understood as
a cache derived from history.

A snapshot can be generated by replaying the entire history,
but systems often store partial summaries in order to
accelerate this reconstruction.

These summaries do not replace history; they merely
compress it.

If the cache becomes inconsistent, the system can
recover by replaying the underlying event log.

\subsection{Causal Correctness}

Replayability ensures that the causal ordering
of events is preserved.

If events are executed out of order or omitted,
the resulting state may diverge from the one
produced by the original history.

Thus correctness is defined not by the final state
alone but by the validity of the event sequence
that generated it.

\begin{principle}[Causal Correctness]
A system state is valid only if it can be reproduced
by replaying a compatible event history.
\end{principle}

\subsection{Failure Recovery}

The replay model also provides a natural mechanism
for recovery.

If a system crashes or loses derived state,
the history remains sufficient to reconstruct
the system.

Distributed systems, transaction logs,
and append-only ledgers already employ this principle.

In these architectures, durability resides
in the event history rather than in the current state.

\subsection{Determinism Without Global Synchronization}

Replayability allows deterministic outcomes
even in systems where events occur concurrently.

If the causal relations between events are preserved,
the final reconstructed state remains stable.

Thus determinism does not require a single global clock
or centralized control. It requires only that the
history preserve causal ordering constraints.

\subsection{Implications}

The event-historical view reframes determinism
as a property of execution history rather than
of instantaneous states.

Deterministic systems therefore become those whose
histories can be faithfully recorded and replayed.

This principle underlies modern architectures such as
event sourcing, distributed logs, and reproducible
build systems.

In the next section we examine how composition
arises from the merging of compatible histories.

% ------------------------------------------------------------
% 5. Composition Through History Merge
% ------------------------------------------------------------

\section{Composition Through History Merge}

Traditional programming models treat composition as the
combination of objects or modules. Systems are assembled
by placing components inside larger containers.

Within an event-historical framework, composition is not
a matter of containment. It is the construction of a
shared history.

Programs compose when their histories are merged.

\subsection{Branching Histories}

Execution frequently produces divergent histories.

Let $H_0$ be a common prefix and suppose two execution
paths extend it:

\[
H_1 = H_0 \cdot A
\qquad
H_2 = H_0 \cdot B
\]

where $A$ and $B$ are sequences of events.

These histories represent two branches of execution.

\subsection{Merge as Joint Extension}

Composition occurs when these branches are reconciled.

\begin{definition}[History Merge]
Given compatible histories $H_1$ and $H_2$ with a common
prefix $H_0$, a merge constructs a history

\[
H_3 = H_0 \cdot A \cup B
\]

that preserves the causal ordering of events in both
branches.
\end{definition}

The resulting history is a joint continuation that
incorporates the contributions of both branches.

\subsection{Causal Preservation}

A valid merge must preserve the causal ordering
within each history.

If event $e_i$ precedes $e_j$ in $H_1$, the same ordering
must hold in the merged history.

\begin{principle}[Causal Preservation]
Merge operations must preserve the partial ordering
of events present in each contributing history.
\end{principle}

This condition ensures that merged histories remain
replayable.

\subsection{Conflict and Compensation}

Not all histories are directly compatible.

If two branches introduce contradictory events,
a merge requires additional events that reconcile
the conflict.

Rather than deleting events, systems introduce
compensating actions that restore consistency.

This preserves the irreversibility of the history.

\subsection{Operational Mereology}

This interpretation provides an alternative to
set-theoretic mereology.

In classical mereology, wholes are defined by
membership relations among parts.

In an operational framework, wholes are defined
by construction history.

A composite system is therefore the result of
merging the histories of its contributing processes.

\subsection{Examples}

The form of composition described by the history algebra
appears in a wide range of computational systems whose
behavior is organized around the accumulation and
reconciliation of event histories.

Distributed version-control systems provide a particularly
clear illustration. In systems such as Git, the evolution
of a repository is represented as a directed acyclic graph
of commits. Each commit corresponds to an event that extends
a prior history, while merge commits reconcile divergent
branches by constructing a new history that incorporates
multiple predecessors. The repository structure therefore
emerges from the composition of histories rather than from
the modification of a single mutable object.

A closely related structure appears in conflict-free
replicated data types (CRDTs), where replicas maintain
operation logs that grow monotonically over time. Updates
are propagated among replicas through merge operations that
combine histories while preserving causal ordering. The
resulting convergence properties arise directly from the
join-semilattice structure imposed on these histories.

Collaborative editing systems exhibit the same pattern at
the level of user interaction. Individual edits are recorded
as operations that extend a shared document history. When
multiple users edit concurrently, their independent
histories are reconciled through merge procedures that
preserve the causal relationships among operations while
producing a coherent continuation of the document.

Event-sourced architectures adopt a similar principle in
the design of large-scale software systems. Instead of
storing mutable state directly, these systems record the
sequence of events that produced that state. The system
state is reconstructed by replaying the event history,
and divergent histories can be reconciled by combining
their logs.

Across these diverse domains, system structure arises
through the composition and reconciliation of histories.
Objects and states appear only as derived reductions of
these histories, while the underlying execution trace
remains the primary carrier of causal structure.

\subsection{Composition Without Global Control}

Because merges operate on histories rather than
global states, composition does not require
central coordination.

Independent processes can extend their local
histories and later reconcile them through
merge operations.

This property allows large systems to scale
without imposing a single controlling authority.

\subsection{Summary}

Within the minimal operational kernel,
composition is defined as the merge of compatible
event histories.

Programs do not combine by nesting objects.
They combine by reconciling the histories
of their execution.

% ------------------------------------------------------------
% 6. Abstraction as Reduction
% ------------------------------------------------------------

\section{Abstraction as Reduction}

If execution constructs histories and composition merges them,
then abstraction must be understood as an operation on those
histories.

Traditional programming theory often treats abstraction
as the introduction of new conceptual entities:
types, classes, modules, interfaces, and other
structural constructs.

Within the event-historical framework, abstraction
is not generative. It is reductive.

Abstraction discards distinctions in a history
while preserving those relevant to a given purpose.

\subsection{Reduction Operators}

\begin{definition}[History Reduction]
A reduction is a mapping
\[
\rho : H \rightarrow H'
\]
that transforms an event history $H$ into a reduced
representation $H'$ by removing information while
preserving selected structural relations.
\end{definition}

Reductions may collapse multiple events into a single
summary, discard ordering relations, or eliminate
distinctions that are irrelevant for a particular task.

\subsection{State as a Reduction}

The most familiar example of reduction is the construction
of system state.

A snapshot summarizes a history by retaining only the
information necessary to describe the system at a
particular moment.

Formally, the state projection
\[
\sigma : H \rightarrow S
\]
is a reduction from the event history into a state space.

This projection discards the sequence of events
that produced the state.

\subsection{Layered Abstractions}

Within the event-historical framework, multiple reductions
may be applied to the same underlying history in order to
derive distinct representations of the system. Each reduction
selects and preserves particular structural properties of
the execution trace while discarding others. The resulting
representations therefore differ in the information they
retain, even though they are derived from the same underlying
history.

For example, a single execution history within a computational
environment may be reduced to produce a filesystem snapshot
that summarizes the persistent storage state of the system.
The same history may also be reduced to generate a database
view reflecting the current contents of structured data
stores, a monitoring representation summarizing resource
utilization and operational metrics, or an audit report
capturing security-relevant transitions and policy events.
Each of these representations is constructed through a
distinct reduction operator applied to the same history.

These derived structures should not be interpreted as
independent states of the system. Instead, they represent
different projections of a common execution trace. Each
projection preserves those aspects of the history relevant
to a particular analytical or operational purpose while
compressing or discarding other details.

From this perspective, the multiplicity of system
representations encountered in practice does not imply
the existence of multiple underlying realities. Rather,
it reflects the existence of multiple reductions applied
to a single evolving history. Layered abstractions therefore
emerge as different informational views of the same
event-historical substrate.

\subsection{Information Loss}

Reduction operators inevitably introduce information loss.
Formally, let
\[
\rho : \mathcal{H} \rightarrow S
\]
be a reduction mapping histories to a space of derived
representations. Because $\rho$ discards structural
information about the internal ordering of events,
it is generally not injective.

Consequently, distinct histories may collapse to the
same reduced representation. That is, there may exist
histories $H_1, H_2 \in \mathcal{H}$ with
\[
H_1 \neq H_2
\]
such that
\[
\rho(H_1) = \rho(H_2).
\]

This property expresses the fundamental loss of
information produced by abstraction.

\begin{principle}[Irreversibility of Reduction]
For histories $H_1, H_2 \in \mathcal{H}$,
\[
\rho(H_1) = \rho(H_2) \;\not\Rightarrow\; H_1 = H_2.
\]
\end{principle}

In other words, reductions cannot generally be inverted.
Once the detailed ordering and structure of a history
have been compressed into a reduced representation,
that information cannot be uniquely recovered from the
reduction alone.

Abstraction therefore functions as a compression
mechanism that maps large execution histories into
representations suitable for reasoning, monitoring,
or control while sacrificing some portion of the
original structural information.

\subsection{Compression and Computation}

The use of reductions serves an important practical
role in computational systems. Execution histories may
grow arbitrarily large as events accumulate over time,
making direct reasoning over the full history expensive
or infeasible. Reduction operators therefore provide
compressed representations that preserve selected
properties of the history while discarding detail
that is irrelevant for a particular computational
task.

These compressed representations enable systems to
perform analysis, coordination, and decision-making
over tractable summaries of their execution. However,
the reduced structures do not replace the underlying
histories. Instead, they function as projections whose
validity ultimately depends upon the histories from
which they were derived.

When ambiguities or inconsistencies arise within a
reduced representation, the system may revert to the
original history and reconstruct the desired state
through replay. In this sense, the history remains the
authoritative source of information, while reductions
provide computationally convenient views that enable
efficient reasoning about the system's behavior.

\subsection{Algorithmic Compression of Histories}

The relationship between abstraction and information loss
can be examined more precisely by considering the
algorithmic compressibility of execution histories.

Let $H \in \mathcal{H}$ denote a finite history consisting
of an ordered sequence of events. A universal interpreter
$U$ may be viewed as a mechanism that reconstructs histories
from descriptions. A program $p$ generates the history $H$
if

\[
U(p) = H.
\]

The algorithmic complexity of the history may then be
defined as

\[
K(H) = \min_{p : U(p) = H} |p|,
\]

where $|p|$ denotes the length of the description $p$.

This quantity measures the shortest description capable
of reconstructing the history through execution.

Reductions can be interpreted as approximations to this
compression process. A reduction operator

\[
\rho : \mathcal{H} \rightarrow S
\]

maps a history to a representation that preserves
certain structural properties while discarding others.
In many cases the reduced representation may be viewed
as a description whose length is significantly smaller
than the full history.

However, unlike the minimal description defined by
algorithmic complexity, reductions are typically
designed to preserve particular observables rather
than to minimize description length globally.

Consequently, the reduced representation should be
understood as a constrained compression of the history
rather than an optimal one.

From this perspective, abstraction functions as a
computationally tractable approximation to the
compression of execution histories. Systems reason
about reduced descriptions of their past behavior,
while the full history remains available for
reconstruction when additional detail is required.

This relationship between histories and their
descriptions further clarifies the role of replay.
Because the history itself constitutes the most
complete description of the system's evolution,
any reduced representation ultimately derives its
validity from the ability to reconstruct behavior
through the replay of the underlying events.

\subsection{Abstraction as Information Geometry}

From a structural perspective, reductions define a
mapping from a high-dimensional history space
to a lower-dimensional representation.

Different abstractions correspond to different
projections of this space.

Thus abstraction can be interpreted as a geometric
operation that reduces dimensionality while preserving
selected invariants.

\subsection{Summary}

Within the minimal operational kernel, abstraction
is understood as the reduction of event histories.

Execution extends histories.

Composition merges histories.

Abstraction reduces histories.

Together these operations provide a unified account
of computation in which state, structure, and meaning
all arise from the manipulation of irreversible
execution histories.

% ------------------------------------------------------------
% 7. The Algebra of Histories
% ------------------------------------------------------------

\section{The Algebra of Histories}

Once execution histories are treated as the primitive
structure of computation, they exhibit a natural algebraic
organization.

This organization arises from the three kernel operations:
extension, merge, and reduction.

\subsection{Prefix Ordering}

Histories are ordered by extension.

\begin{definition}[Prefix Order]
Given two histories $H_1$ and $H_2$, we say that
\[
H_1 \preceq H_2
\]
if $H_1$ is a prefix of $H_2$.
\end{definition}

This relation defines a partial order over the set
of all histories.

\begin{proposition}
The set of event histories equipped with the prefix
relation forms a partially ordered set.
\end{proposition}

Intuitively, one history precedes another when it
represents an earlier stage of execution.

\subsection{Monotonic Growth}

Execution always moves forward in this ordering.

\begin{principle}[Monotonic Execution]
If an event $e$ extends history $H$ to $H'$,
then
\[
H \preceq H'.
\]
\end{principle}

This property ensures that computation cannot
erase history; it can only extend it.

\subsection{Merge Semilattice}

Compatible histories admit a merge operation.

\begin{definition}[History Join]
Given compatible histories $H_1$ and $H_2$,
their merge is the least history $H_3$ such that
\[
H_1 \preceq H_3
\qquad
H_2 \preceq H_3.
\]
\end{definition}

This operation corresponds to the join in a
semilattice.

\begin{proposition}
The set of compatible histories equipped with
the merge operator forms a join-semilattice.
\end{proposition}

The merged history represents the smallest
history containing both branches.

\subsection{Concurrency}

Because histories are partially ordered rather
than strictly sequential, independent events
may occur in different orders without affecting
the resulting state projection.

Such events are concurrent.

\begin{definition}[Concurrent Events]
Events $e_i$ and $e_j$ are concurrent if neither
\[
e_i \leq e_j
\quad \text{nor} \quad
e_j \leq e_i.
\]
\end{definition}

Concurrency therefore corresponds to independence
in the causal ordering.

\subsection{Reduction Morphisms}

Reductions act as structure-preserving mappings
between histories.

\begin{definition}[Reduction Morphism]
A reduction $\rho$ is a mapping between histories
that preserves causal relations among retained
events.
\end{definition}

These mappings compress histories while maintaining
the ordering relations necessary for replay.

\subsection{Derived Structure}

The preceding analysis implies that the space of histories
admits a well-defined algebraic structure. Histories are
organized by a prefix partial order in which one history
is considered less than or equal to another when it
constitutes an initial segment of that history. This
ordering relation captures the causal containment of
execution traces and provides the fundamental temporal
structure of the system.

Within this ordered space, execution corresponds to a
monotonic extension operator that appends new events
while preserving the ordering constraints of prior
transitions. Histories therefore evolve through the
irreversible accumulation of events whose effects
propagate through the existing causal structure.
When independent histories are compatible, they may
be combined through a join-like operation that
constructs a minimal history containing both as
substructures while preserving their causal ordering.

In addition to these growth operations, the algebra
admits reduction morphisms that map histories to
compressed representations such as states, snapshots,
or higher-level abstractions. These reductions discard
information about intermediate transitions while
preserving selected structural properties of the
underlying execution trace.

Taken together, these components define a computational
structure whose behavior closely resembles that of
systems already encountered in practice, including
distributed append-only logs, version-control graphs,
and event-sourced architectures. In each of these
settings, system evolution is governed by the monotonic
extension of histories, the reconciliation of compatible
branches, and the projection of histories into derived
representations.

\subsection{Interpretation}

The algebraic structure described above demonstrates
that the event-historical kernel is not merely a
conceptual reformulation of computation. Instead, it
corresponds directly to mathematical structures that
already appear in the design of large-scale computational
systems.

Modern distributed infrastructures routinely employ
append-only logs, causal histories, and merge-based
reconciliation mechanisms in order to maintain
consistency across independent replicas. Similarly,
version-control systems encode the evolution of
software artifacts as directed graphs of commits,
while event-sourced architectures treat persistent
event histories as the authoritative record from
which system state is derived.

What distinguishes the present framework is not the
existence of these structures themselves, but the
interpretive shift in which they are regarded as
foundational rather than incidental. In conventional
models, such mechanisms are typically understood as
implementation techniques layered upon an underlying
state-based semantics. Within the event-historical
kernel, by contrast, histories are treated as the
primary objects of computation, while states and
abstractions arise only as reductions of those
histories.

This perspective therefore repositions a class of
practical engineering patterns as manifestations of
a more general algebra of execution, suggesting that
the operational behavior of large-scale computational
systems reflects deeper structural principles of
history-based computation.

% ------------------------------------------------------------
% Formalization of the History Algebra
% ------------------------------------------------------------

\subsection{Formalization of the History Algebra}

The structures described informally above can be expressed
more precisely by introducing a minimal algebra of histories.
Let $\mathcal{H}$ denote the set of all admissible histories
generated by a computational system.

\subsubsection{History Poset}

Histories are ordered by a prefix relation capturing causal
containment. For histories $H_1, H_2 \in \mathcal{H}$ we write

\[
H_1 \leq H_2
\]

when $H_1$ is a prefix of $H_2$, meaning that the sequence
of events constituting $H_1$ appears as an initial segment
of the events of $H_2$ while preserving causal ordering.

The pair $(\mathcal{H}, \leq)$ therefore forms a partially
ordered set in which earlier histories represent causal
substructures of later ones.

\subsubsection{Extension Operator}

Execution corresponds to the monotonic extension of histories.
Let $E$ denote the set of admissible events. The extension
operator is a partial function

\[
\operatorname{ext} : \mathcal{H} \times E \rightarrow \mathcal{H}
\]

such that if $H' = \operatorname{ext}(H, e)$ then

\[
H \leq H'.
\]

This operator appends event $e$ to history $H$ while
preserving the causal constraints required for the event
to occur. Because extension always produces a larger
history, execution is monotonic with respect to the
prefix ordering.

\subsubsection{Merge and Join}

When two histories are compatible—meaning that their
events do not violate causal constraints—they may be
combined through a merge operation.

Formally, if $H_1$ and $H_2$ are compatible histories,
their merge is defined as a least upper bound

\[
H_1 \sqcup H_2
\]

satisfying

\[
H_1 \leq H_1 \sqcup H_2, \qquad
H_2 \leq H_1 \sqcup H_2,
\]

and for any history $H$ such that
$H_1 \leq H$ and $H_2 \leq H$,

\[
H_1 \sqcup H_2 \leq H.
\]

The merge therefore acts as a join operation over
compatible histories. When defined for all compatible
pairs, the space of histories forms a join-semilattice.

\subsubsection{Reduction Morphisms}

Many computational systems maintain derived representations
of histories in the form of states, snapshots, or summaries.
These representations can be modeled as reduction morphisms

\[
\sigma : \mathcal{H} \rightarrow S,
\]

where $S$ is a space of reduced representations.

A reduction discards some information about the internal
structure of a history while preserving selected properties.
Multiple histories may therefore map to the same reduced
state:

\[
\sigma(H_1) = \sigma(H_2)
\]

even when $H_1 \neq H_2$.

\subsubsection{Summary}

The resulting algebra consists of a history poset
$(\mathcal{H}, \leq)$ equipped with a monotone extension
operator, a join operation over compatible histories,
and reduction morphisms mapping histories to derived
representations. These structures together provide a
minimal formal description of computation as the
irreversible construction and transformation of
event histories.

% ------------------------------------------------------------
% Fundamental Properties of the History Algebra
% ------------------------------------------------------------

\subsection{Fundamental Properties}

The algebra of histories admits several structural
properties that follow directly from the definitions
introduced above. These properties clarify how
determinism, convergence, and reconstruction arise
within the event-historical kernel.

\begin{theorem}[Deterministic Replay]
Let $H \in \mathcal{H}$ be a history and let
$\sigma : \mathcal{H} \rightarrow S$ be a reduction
operator that reconstructs a state representation
from the history. If the event semantics are
deterministic, then replaying the events of $H$
produces a unique state $\sigma(H)$.
\end{theorem}

\begin{proof}
Replay reconstructs the system state by applying
the events of $H$ in causal order starting from an
initial configuration. Because the semantics of
each event are deterministic and because the causal
ordering of $H$ is preserved by the prefix relation,
each event produces a uniquely determined
transformation of the intermediate configuration.
The sequential application of these transformations
therefore produces a unique final configuration.
Consequently the mapping $\sigma(H)$ is well-defined
and deterministic.
\end{proof}

\begin{theorem}[Monotonic Extension]
For any history $H$ and admissible event $e$ such that
$H' = \operatorname{ext}(H,e)$ is defined, the prefix
ordering satisfies
\[
H \leq H'.
\]
\end{theorem}

\begin{proof}
By definition, the extension operator appends event $e$
to the existing sequence of events in $H$ while preserving
their causal ordering. The resulting history therefore
contains $H$ as a prefix. Hence $H \leq H'$.
\end{proof}

\begin{theorem}[Merge Convergence]
Let $H_1$ and $H_2$ be compatible histories.
If the join operation $H_1 \sqcup H_2$ exists,
then the merged history is the least upper bound
of the pair under the prefix ordering.
\end{theorem}

\begin{proof}
Compatibility implies that the events of $H_1$ and
$H_2$ do not violate each other's causal constraints.
The merge operation constructs a history containing
both histories as prefixes while preserving causal
ordering. Any history containing both $H_1$ and $H_2$
must contain at least the events present in their
merge. Therefore the merge constitutes the least
upper bound under the prefix order.
\end{proof}

\subsection{Interpretation}

These properties illustrate that deterministic replay,
monotonic execution, and merge-based convergence are
not independent engineering techniques but structural
consequences of the history algebra itself. Systems
built upon event logs, distributed histories, and
append-only structures therefore inherit these
properties directly from the algebraic organization
of their execution traces.

The operational kernel described in this work can
thus be understood as a minimal formal structure
in which computation is modeled as the irreversible
construction and reconciliation of histories subject
to causal constraints.

% ------------------------------------------------------------
% 8. Time, Causality, and Irreversibility
% ------------------------------------------------------------

\section{Time, Causality, and Irreversibility}

If computation is defined as the irreversible extension of
event histories, then temporal structure arises naturally
from the ordering relations within those histories.

Rather than treating time as an independent dimension,
the framework interprets time as a property of execution.

\subsection{Time as History Length}

The simplest notion of time corresponds to the growth
of a history.

\begin{definition}[Execution Time]
The execution time of a history $H$ is the length of
its event sequence.
\[
t(H) = |H|
\]
\end{definition}

Each extension of the history advances time by one step.

In partially ordered histories, time corresponds to the
depth of the causal structure rather than a global clock.

\subsection{Causality}

Causality emerges from the ordering relation
defined over events.

\begin{definition}[Causal Relation]
Event $e_i$ causally precedes event $e_j$ if
\[
e_i \leq e_j
\]
within the history poset.
\end{definition}

This ordering captures the dependency relations
that govern execution.

If two events are incomparable in the ordering,
they may occur independently.

\subsection{Irreversibility}

The extension operator is inherently irreversible.

Once an event has been appended to a history,
it cannot be removed without introducing additional
events that compensate for its effects.

\begin{principle}[Historical Irreversibility]
Execution histories grow monotonically and cannot
be reduced without loss of information.
\end{principle}

This property distinguishes event-historical systems
from reversible state-transition models.

\subsection{Entropy and Information Loss}

Reductions introduce a second and distinct form of irreversibility
within the event-historical framework. Whereas execution produces
irreversibility through the monotonic accumulation of events,
abstraction produces irreversibility through the loss of
information about the history from which a representation was
derived.

When a history is reduced to a state, snapshot, or higher-level
abstraction, the resulting representation preserves only a
subset of the structural information contained in the original
execution trace. The ordering of events, the intermediate
configurations encountered during execution, and the causal
structure that produced the final configuration may all be
partially or completely discarded. Multiple distinct histories
may therefore collapse to the same reduced representation.

Two forms of irreversibility therefore coexist within the
operational kernel. The first arises from execution itself,
in which events accumulate irreversibly to extend a history.
The second arises from abstraction, in which reductions compress
histories by discarding information about the sequence of
transitions that produced a configuration. These two processes
operate simultaneously: histories grow through execution while
representations of those histories become progressively more
coarse through reduction.

This dual structure closely resembles the relationship between
microscopic dynamics and macroscopic description in thermodynamic
systems. At the microscopic level, transitions accumulate through
local interactions that preserve detailed causal structure. At
the macroscopic level, thermodynamic variables summarize the
system by discarding the majority of that detail. The increase
of entropy in such systems reflects the growth of possible
microscopic histories consistent with a given macroscopic state.

Within the event-historical kernel, reductions play an analogous
role. As histories are compressed into states or abstractions,
the number of distinct histories compatible with a given
representation increases. The loss of information produced by
reduction therefore functions as a computational analogue of
entropy, reflecting the growing indeterminacy of the underlying
execution histories once their detailed structure has been
discarded.

\subsection{The Arrow of Execution}

Within the event-historical framework, the direction of time
is naturally identified with the direction in which histories
extend. A computation is not a traversal through a fixed
space of states but the progressive construction of an
execution history through the accumulation of events.
Each event appends new structure to the history while
preserving the causal ordering of prior transitions.

This monotonic extension induces a natural temporal
orientation. If $H_1 \leq H_2$ denotes that history
$H_2$ extends $H_1$ by the addition of further events,
then the ordering relation defines a partial temporal
structure over histories. Valid computations therefore
proceed only in the direction of increasing histories.
The arrow of execution corresponds to this monotone
growth of causal structure.

Attempts to ``reverse'' execution do not actually undo
events within the system. Instead, they require selecting
an earlier prefix of the history and re-executing the
subsequent computation from that point. The underlying
history itself remains irreversible; what appears as
reversal is merely the construction of a new extension
that begins from a previously recorded prefix.

\subsection{Replay and Reconstruction}

Because execution histories are preserved, past states
of the system may be reconstructed through replay.
Given a history $H$ and a projection operator
$\sigma(H)$ that derives a state representation
from that history, replay corresponds to the
reconstruction of $\sigma(H)$ by sequentially
reapplying the events recorded in the history.

Importantly, this process does not reverse the
system’s evolution. Replay merely reconstructs
earlier configurations that were produced during
the growth of the history. The system itself
continues to evolve through the monotonic extension
of its execution trace.

The relationship between replay, reduction, and
information loss may be expressed more precisely
through an entropy-like quantity defined over
reduced representations of histories. If
$\sigma(H)$ denotes a reduction of the history
$H$ into a state representation, the number of
distinct histories compatible with that reduction
may be written as

\[
\Omega(\sigma(H)) =
\left| \{ H' \mid \sigma(H') = \sigma(H) \} \right|.
\]

An entropy associated with the reduction can then
be defined as

\[
S(\sigma(H)) = \log \Omega(\sigma(H)).
\]

This quantity measures the multiplicity of execution
histories that collapse to the same reduced state.
As reductions discard information about causal
ordering and intermediate transitions, the number
of histories consistent with a given representation
increases. The entropy therefore reflects the
degree of indeterminacy introduced by abstraction.

Within this formulation, execution irreversibility
and abstraction irreversibility interact in a manner
analogous to microscopic and macroscopic dynamics in
thermodynamic systems. Histories grow through the
irreversible accumulation of events, while reductions
coarse-grain those histories into representations that
discard detail. Replay allows the reconstruction of
earlier configurations, but the underlying history
continues to extend forward, preserving the arrow
of execution.

\subsection{Interpretation}

The event-historical kernel developed in the preceding
sections provides a unified interpretation of temporal
structure within computational systems. Rather than
treating time, causality, and irreversibility as
independent primitives that must be introduced through
external assumptions, the kernel demonstrates that
these properties arise naturally from the algebraic
structure of execution histories.

Within this framework, time is not modeled as an
independent dimension in which computation unfolds.
Instead, temporal structure emerges from the monotonic
extension of histories through the accumulation of
events. The ordering relation defined over histories
induces a partial temporal structure in which earlier
events are those contained within the prefixes of
later histories. The progression of time therefore
corresponds directly to the growth of execution
histories themselves.

Causal structure likewise arises from the ordering
relations that constrain the extension of histories.
Events are not merely arranged sequentially but are
linked by relations that record the dependencies
required for their occurrence. The resulting partial
order captures the propagation of constraints through
the system and provides the foundation for deterministic
replay and reproducibility. In this sense, causality is
not an external principle governing computation but a
structural property of the histories through which
computation is realized.

Irreversibility follows from two distinct but related
features of the kernel. The first arises from the
monotonic growth of histories: once an event has been
incorporated into a history, it cannot be removed
without constructing a new history that explicitly
records the compensating transition. The second arises
from reductions that compress histories into states or
abstractions by discarding information about the
intermediate transitions that produced them. These
reductions introduce an additional form of irreversibility
by collapsing multiple distinct histories into a single
representation.

Taken together, these properties demonstrate that the
temporal asymmetries observed in computational and
physical systems do not require independent axiomatic
assumptions. They emerge directly from the algebra of
histories defined by the event-historical kernel. Time,
causality, and irreversibility therefore appear not as
external constraints imposed upon computation, but as
structural consequences of the irreversible construction
process through which systems evolve.

\subsection{Event Histories and Local Constraint Systems}

The event-historical kernel described above is not merely
a computational abstraction. Systems with similar structure
appear in statistical physics, particularly in lattice models
such as the Ising system.

In the Ising model, a lattice is composed of sites indexed by
$i$, each of which carries a local spin variable

\[
s_i \in \{-1, +1\}.
\]

The energy of the lattice configuration is determined by a
Hamiltonian of the form

\[
E = -J \sum_{\langle i,j \rangle} s_i s_j - h \sum_i s_i ,
\]

where $J$ represents the coupling between neighboring sites,
$h$ represents an external field, and $\langle i,j \rangle$
denotes neighboring lattice pairs.

The configuration of the lattice evolves through local spin
updates. Each update changes the spin at a site and thereby
modifies the interaction energies with neighboring sites.

\subsubsection{Spin Flips as Events}

Within the event-historical framework, each spin flip can be
interpreted as an event

\[
e_k = (i_k, \Delta s_{i_k}),
\]

where $i_k$ denotes the lattice site and $\Delta s_{i_k}$ is
the change in the local spin.

An execution history of the lattice therefore takes the form

\[
H = (e_1, e_2, \dots, e_n),
\]

representing the sequence of spin updates that produced the
current configuration.

\subsubsection{State as Derived Configuration}

The instantaneous lattice configuration can be recovered by
replaying this sequence of events from an initial configuration.

Thus the spin state vector

\[
\mathbf{s} = (s_1, s_2, \dots, s_N)
\]

is a derived state obtained through the projection

\[
\sigma(H) = \mathbf{s}.
\]

The configuration itself therefore does not contain the
complete causal information of the system; it is a reduction
of the underlying event history.

\subsubsection{Constraint Propagation}

Each spin flip modifies the local energy landscape of the
system. When a site changes state, the interaction energies
with its neighbors are altered.

This change constrains which future spin flips are energetically
favorable or even possible.

Execution therefore proceeds through a sequence of events in
which each event reshapes the constraints governing subsequent
events.

\subsubsection{Material Reordering}

From this perspective, the lattice may be interpreted as a
physical substrate whose set of available transitions is
continually reorganized by prior actions. Rather than
representing a static collection of possible states, the
structure evolves as events modify the local configuration
of relations among components. Each event therefore alters
the material conditions under which subsequent transitions
may occur, effectively reshaping the landscape of allowable
operations.

In such systems, execution is inseparable from the process
by which the system reorders the resources available to it.
Local events do not merely advance a computation; they also
transform the constraints that govern future activity.
The causal ordering of events thus reflects the propagation
of interaction constraints across the structure of the
system. What appears externally as a sequence of operations
is internally a progressive reconfiguration of the relations
that define the system’s permissible dynamics.

This interpretation aligns directly with the event-historical
kernel developed earlier. Execution proceeds through the
irreversible extension of histories by locally constrained
events whose ordering records the propagation of interaction
constraints. The configurations that appear as observable
states of the system are therefore reductions of the
underlying event history, summarizing the cumulative
reordering produced by prior transitions.

Under this interpretation, computation is not the traversal
of a pre-existing space of states but a process of
material reorganization in which each event reshapes the
structure of possibilities available to the system.

\subsubsection{Interpretation}

The Ising system therefore provides a concrete physical
analogue of the kernel structure developed earlier.

Rather than evolving through abstract state transitions,
the lattice evolves through the irreversible accumulation
of local events whose effects propagate through neighboring
constraints.

The instantaneous configuration of the lattice is thus
best understood as a compressed description of its
execution history.

This correspondence suggests that the event-historical
kernel reflects a structural principle shared by both
computational and physical systems: dynamics arise from
the monotonic accumulation of locally constrained events.

\subsection{Monotone Constraint Dynamics}

The correspondence between event histories and lattice dynamics
can be made more precise by observing that both systems evolve
through monotone constraint updates.

Let $H$ denote an execution history and let $\mathcal{C}(H)$
denote the set of constraints induced by that history.
Each event modifies this constraint structure.

Formally, if an event $e$ extends the history,

\[
H' = H \cdot e ,
\]

then the associated constraint set evolves according to

\[
\mathcal{C}(H') = F(\mathcal{C}(H), e),
\]

where $F$ updates the constraints generated by the system.

\subsubsection{Constraint Monotonicity}

In many systems the constraint structure evolves monotonically.

\begin{principle}[Constraint Monotonicity]
For histories $H \preceq H'$,

\[
\mathcal{C}(H) \subseteq \mathcal{C}(H').
\]
\end{principle}

In other words, execution does not remove previously
established constraints; it only introduces new ones
or strengthens existing ones.

\subsubsection{Energy Interpretation}

This structure admits an energy-like interpretation
analogous to that used in statistical physics.

Let

\[
E(H)
\]

denote a scalar potential defined over histories.
Execution tends to move the system toward configurations
that reduce this potential.

Each event therefore produces a change

\[
\Delta E = E(H') - E(H).
\]

In Ising dynamics, spin flips are accepted when they
reduce the system's energy or when thermal fluctuations
permit a temporary increase.

Similarly, execution histories evolve through events
that locally reorganize constraints in ways that tend
toward stable configurations.

\subsubsection{Local Update Rule}

In lattice systems the update rule is local.
A spin flip at site $i$ affects only neighboring
sites.

This locality corresponds to the constraint update
function $F$ operating on a bounded region of the
system.

Thus the constraint evolution

\[
\mathcal{C}(H') = F(\mathcal{C}(H), e)
\]

depends only on the local neighborhood of the event.

\subsubsection{Descent Dynamics}

Under repeated execution, the system performs a form
of constrained descent over the potential landscape.

The history therefore traces a path through a
configuration manifold governed by local interactions.

This process is analogous to the relaxation dynamics
of spin systems, where local updates gradually
reduce the system's global energy.
\subsubsection{Structural Equivalence}

The comparison developed in the preceding sections
reveals a structural correspondence between the
event-historical kernel and classes of physical
systems governed by local constraint dynamics.
Although these domains are typically analyzed using
different mathematical frameworks, their underlying
organizational principles exhibit a striking degree
of similarity when expressed in terms of histories
and constraint propagation.

In both cases, system evolution proceeds through
local transitions that alter the configuration of
neighboring components while preserving the causal
ordering of prior interactions. Each transition
introduces additional constraints into the system,
and these constraints accumulate monotonically as
execution progresses. The resulting dynamics are not
arbitrary but are guided toward configurations that
minimize inconsistencies among interacting elements,
producing a gradual descent toward stable or
metastable arrangements.

Within the event-historical kernel, this process
appears as the irreversible extension of histories
through compatible events whose effects propagate
across the structure of the system. In physical
lattice models, an analogous process occurs through
sequences of local state updates—such as spin flips
in Ising-type systems—that progressively reorganize
the configuration of the lattice under the influence
of interaction potentials.

From this perspective, the event history of a
computational system plays a role closely analogous
to the sequence of local updates that drive the
dynamics of a constrained physical lattice. In both
cases, the observable configuration at any moment
can be interpreted as a reduction of the cumulative
history of transitions that produced it. Execution
therefore corresponds to a trajectory through a
constrained energy landscape defined by the
interactions among local components, with stability
emerging as the system approaches configurations
that satisfy the constraints imposed by its
structure.

\subsubsection{Interpretation}

This correspondence suggests that the kernel
describes not only a model of computation but
a broader class of dynamical systems.

In both computational and physical settings,
structure emerges from the cumulative effects
of locally constrained transitions.

Histories record those transitions,
and the resulting states represent compressed
descriptions of the underlying execution path.

\subsection{History Lattices and Potential Functions}

The algebra introduced earlier can be refined further by
observing that the space of histories forms a lattice-like
structure equipped with a monotone potential.

Each history represents a partially ordered sequence of
events whose causal relationships constrain the set of
possible extensions. When two compatible histories are
combined through composition, the resulting structure
preserves the ordering constraints of both inputs while
producing a common continuation. This behavior resembles
the join operation of a semilattice, in which compatible
elements admit a least upper bound representing their
combined information.

Within this framework, histories can therefore be viewed
as elements of a partially ordered space in which extension
corresponds to monotonic growth. If $H_1 \leq H_2$ denotes
that history $H_2$ extends $H_1$ by the addition of further
events, then the ordering relation captures the causal
containment of execution traces. The merge of compatible
histories corresponds to the construction of a minimal
history that contains both as substructures while preserving
their causal ordering.

This ordering admits the introduction of a potential
function defined over histories. The potential measures
the degree to which a history satisfies the constraints
governing the system. Each event alters the potential by
modifying the local configuration of constraints, and the
evolution of execution can therefore be interpreted as a
trajectory through the space of histories guided by this
potential landscape.

In many systems, the potential decreases as constraints
become satisfied or inconsistencies are eliminated.
Execution thus exhibits a form of constrained descent
toward stable configurations in which further transitions
no longer reduce the potential. These stable histories
correspond to fixed points of the execution dynamics and
appear externally as equilibrium states of the system.

The introduction of a potential over the history lattice
provides a bridge between computational execution and
physical models governed by local interaction energies.
Just as spin systems evolve through local updates that
reduce an energy functional, event-driven systems evolve
through transitions that reorganize the constraints of
their histories. In both cases, the global structure of
the system emerges from the cumulative effect of locally
constrained transitions.

\subsubsection{Join-Semilattice Structure}

Let $\mathcal{H}$ denote the set of all valid event histories.

The prefix ordering relation

\[
H_1 \preceq H_2
\]

induces a partial order over $\mathcal{H}$.

When compatible histories admit a merge operation,
this structure becomes a join-semilattice.

\begin{definition}[History Join]
For compatible histories $H_1$ and $H_2$, the join
\[
H_1 \vee H_2
\]
is the least history containing both histories as prefixes.
\end{definition}

\begin{proposition}
If merge operations are associative, commutative,
and idempotent, then $(\mathcal{H}, \vee)$ forms
a join-semilattice.
\end{proposition}

These properties appear naturally in systems where
histories accumulate monotonically and merges
preserve causal structure.

\subsubsection{Monotone Potential}

We now introduce a scalar function

\[
E : \mathcal{H} \rightarrow \mathbb{R}
\]

which assigns a potential value to each history.

This potential measures the degree to which
the current configuration satisfies the system's
constraints.

\begin{definition}[Monotone Potential]
A potential function $E$ is monotone if
\[
H_1 \preceq H_2 \Rightarrow E(H_2) \leq E(H_1).
\]
\end{definition}

Execution therefore tends to move the system toward
histories with lower potential.

\subsubsection{Local Descent}

If events modify the system only locally,
then the change in potential depends only
on a bounded region of the history.

\[
\Delta E = E(H \cdot e) - E(H)
\]

is therefore determined by the local interaction
structure associated with the event $e$.

This condition mirrors the local update rules
of lattice systems such as the Ising model.

\subsubsection{Energy Landscape Interpretation}

The set of histories can therefore be interpreted
as a configuration manifold over which the potential
function $E$ defines an energy landscape.

Execution corresponds to a trajectory through this
landscape generated by successive events.

Local updates modify the landscape by changing
the constraints that govern future events.

\subsubsection{Stability and Fixed Points}

Stable configurations arise when further events
cannot reduce the system's potential.

\begin{definition}[Stable History]
A history $H$ is stable if no valid extension
\[
H' = H \cdot e
\]
exists such that
\[
E(H') < E(H).
\]
\end{definition}

Such histories correspond to fixed points of the
execution dynamics.

\subsubsection{Connection to Distributed Systems}

This structure closely resembles the semantics of
conflict-free replicated data types (CRDTs),
distributed logs, and version-control graphs.

In those systems, histories grow monotonically,
merges form semilattice joins, and convergence
arises from the algebraic properties of the merge
operation.

\subsection{Dual Divisibility and the Direction of Execution}

A subtle structural feature of history-based computation arises once one
asks whether every history admits a lawful factorization through intermediate
cuts. In ordinary state-transition models this question rarely appears,
because transitions are defined directly between states. In a history
framework, however, the question becomes fundamental.

Let $\mathcal{H}$ denote the category of executable histories. Objects
represent interfaces or system boundaries, and morphisms represent
history segments. Composition
\[
h_{u\to t} = h_{s\to t} \circ h_{u\to s}
\]
corresponds to temporal concatenation of histories.

The extension operation introduced earlier corresponds precisely to
composition in $\mathcal{H}$. A history prefix is extended by composing
with a subsequent segment.

We now ask whether every history can be decomposed at arbitrary
intermediate cuts.

\begin{definition}[Left divisibility]
A history $h_{0\to t}$ is \emph{left divisible} if for every cut
$0 \le s \le t$ there exists an admissible history segment
$k_{s\to t}$ such that
\[
h_{0\to t} = k_{s\to t} \circ h_{0\to s}.
\]
\end{definition}

Intuitively, left divisibility states that a global execution can be
reconstructed by extending any valid prefix with a lawful continuation.
This corresponds directly to the operational meaning of extension in
the minimal kernel.

However, there exists a second notion of factorization.

\begin{definition}[Right divisibility]
A history $h_{0\to t}$ is \emph{right divisible} if for every cut
$0 \le s \le t$ there exists an admissible history segment
$r_{0\to s}$ such that
\[
h_{0\to t} = h_{s\to t} \circ r_{0\to s}.
\]
\end{definition}

Although these definitions appear symmetric, they need not coincide.
The distinction becomes visible whenever admissibility constraints are
defined relative to an ordered structure over histories.

\begin{theorem}[Dual divisibility of operational representations]
Let histories be represented operationally in two dual ways:
\[
F : \mathcal{H} \to \mathbf{OrdMon}, \qquad
G : \mathcal{H}^{op} \to \mathbf{OrdMon},
\]
where $F$ acts on states and $G$ acts on effects. Suppose the two
representations are prediction-equivalent in the sense that
\[
\langle F(h)(\rho),A\rangle = \langle \rho,G(h)(A)\rangle
\]
for all states $\rho$, effects $A$, and histories $h$.

Then left divisibility of histories relative to the state order does
not in general imply right divisibility relative to the effect order.
\end{theorem}

This asymmetry arises because the admissibility cones governing
state transformations and observable transformations differ.
Operationally, the distinction corresponds to two different ways of
generating histories in continuous time. One obtains either
\[
\dot{\Phi}_t = L_t \circ \Phi_t
\qquad \text{or} \qquad
\dot{\Phi}_t = \Phi_t \circ R_t,
\]
corresponding respectively to postcomposition and precomposition
of infinitesimal execution steps.

Consequently, execution histories possess an intrinsic directional
structure. Extension corresponds to lawful postcomposition of
history segments, while the dual observable representation induces
a precomposition dynamics. These two execution monoids are
prediction-equivalent but structurally distinct.

In the language of the minimal operational kernel, this observation
reveals that the extension operation is fundamentally asymmetric.
The admissibility of intermediate execution steps depends on the
direction in which history factorization is evaluated. Merge and
reduction therefore operate over histories whose compositional
structure is not strictly self-dual.

\subsection{History Cuts, Merge Consistency, and Reduction}

The divisibility properties introduced above have a direct operational
interpretation in the history-based model of computation developed in
this paper. In particular, they determine whether histories can be
safely partitioned, merged, and reduced without violating the
admissibility constraints of the execution kernel.

Consider a history segment
\[
h_{0\to t}.
\]
A cut at time $s$ partitions this history into two segments
\[
h_{0\to s}, \qquad h_{s\to t}.
\]

If the history is left divisible, there exists an admissible
continuation map $k_{s\to t}$ such that
\[
h_{0\to t} = k_{s\to t} \circ h_{0\to s}.
\]

Operationally, this guarantees that any valid prefix can be extended
without ambiguity. In other words, the prefix contains sufficient
information for lawful continuation of the computation.

Merge operations impose a stronger requirement. Suppose two histories
share a common prefix
\[
h_{0\to s}.
\]

We may attempt to merge their continuations,
\[
h_{0\to t_1}, \qquad h_{0\to t_2},
\]
into a combined execution structure. Such a merge is admissible only
when the intermediate extensions remain compatible within the ordered
structure of histories.

Divisibility therefore acts as a consistency condition for merge.
If intermediate extensions fail to remain admissible at every cut,
histories cannot be safely merged without introducing violations of
the execution constraints.

Reduction operations introduce a complementary phenomenon. A reduction
map
\[
R : \mathcal{H} \to \widetilde{\mathcal{H}}
\]
removes distinctions between histories by collapsing them into
coarser equivalence classes. Typical examples include snapshot
generation, aggregation, and state summarization.

For reductions to preserve executability, they must respect the
factorization structure of histories. That is, whenever a history
admits a valid decomposition,
\[
h_{0\to t} = h_{s\to t} \circ h_{0\to s},
\]
the reduced histories must admit a compatible decomposition
\[
R(h_{0\to t}) = R(h_{s\to t}) \circ R(h_{0\to s}).
\]

If this property fails, reduction destroys information necessary for
lawful extension and merge operations.

Consequently, divisibility can be interpreted as a structural
compatibility condition for the minimal operational kernel.
Extension guarantees the existence of lawful continuations,
merge guarantees compatibility between concurrent histories,
and reduction guarantees that coarse-grained representations
preserve the compositional structure of execution.

Together these conditions define the domain in which history-based
computation remains stable under composition and abstraction.

\subsubsection{Interpretation}

The event-historical kernel therefore admits a
precise mathematical interpretation.

Histories form a partially ordered semilattice,
execution corresponds to monotone extension,
and system dynamics can be interpreted as descent
over a constraint-induced potential landscape.

This structure appears in both computational
architectures and physical systems governed by
local interactions.

% ------------------------------------------------------------
% 9. Unified Constraint-Driven Computation
% ------------------------------------------------------------

\section{Unified Constraint-Driven Computation}

The minimal operational kernel introduced in this essay
is not confined to a particular computational paradigm.
Rather, the structural principles that arise from the
event-historical framework appear across a wide range of
computational architectures and physical dynamical systems.
Although these systems are typically studied within
separate theoretical traditions, they exhibit a shared
organizational pattern when interpreted through the lens
of event histories and constraint propagation.

Within this perspective, execution is understood as the
monotonic extension of histories through the occurrence
of events whose effects accumulate irreversibly over time.
Each event modifies the configuration of the system in a
manner that reshapes the constraints governing subsequent
transitions. Because these constraints arise through
local interactions among components, the resulting
dynamics propagate through the system incrementally
rather than through centralized control. Divergent
histories produced by independent activity may later be
reconciled through merge operations that preserve the
causal ordering of events while combining their effects
into a common continuation.

As these locally constrained transitions accumulate,
the system evolves toward configurations that satisfy
the interaction constraints imposed by its structure.
Stable states therefore emerge not as primary objects
of computation but as reductions of the underlying
execution history. From this standpoint, apparently
distinct systems—ranging from distributed computing
architectures to physical lattice models governed by
local interactions—can be interpreted as instances of
a broader class of constraint-driven processes whose
behavior is determined by the irreversible accumulation
of events and the propagation of constraints across
their histories.

\subsection{Deterministic Replay Systems}

Event-sourced computational architectures provide a clear
illustration of the event-historical semantics developed
in this essay. In such systems, the primary persistent
structure is not a mutable state but an append-only log
that records the sequence of events through which the
system has evolved. Each entry in the log represents a
discrete operation that extends the history of the system,
and the cumulative sequence of these operations constitutes
the authoritative record of its execution.

Within this architectural model, the observable state of
the system is not stored as a primary object but instead
is reconstructed by replaying the event log. The sequence
of recorded events therefore corresponds directly to the
history $H$ described in the operational kernel. The
reconstruction process itself is an instance of the
projection mapping $\sigma(H)$ that derives a state
representation from the underlying history. Snapshots
or checkpoints that are sometimes stored to accelerate
recovery can be interpreted as reductions of this history,
compressing the sequence of events into a summary that
preserves the current configuration while discarding
details of the intermediate transitions.

Deterministic replay thus arises naturally within the
event-historical framework. Because the system state is
defined as a projection of the recorded history, identical
histories necessarily produce identical derived states
when replayed. The append-only structure of the event log
ensures that execution histories grow monotonically, while
the replay mechanism guarantees that the causal structure
of execution remains recoverable even after failures or
system restarts.

\subsection{Distributed Convergence}

Distributed systems frequently rely on structural
mechanisms that ensure convergence among replicas
without requiring centralized coordination or a
globally synchronized clock. A well-known example
is the class of conflict-free replicated data types
(CRDTs), which provide deterministic convergence
across distributed nodes despite the presence of
asynchronous communication and independent local
updates.

In such systems, each replica maintains a history
of operations whose structure grows monotonically
over time. Local updates extend this history,
while synchronization between replicas occurs
through merge operations that reconcile divergent
branches of execution. These merges behave
algebraically as join operations over a semilattice,
ensuring that the combined history preserves the
information contained in each contributing branch.
Because the merge operation is associative,
commutative, and idempotent, repeated synchronization
among replicas leads to eventual convergence even
in the absence of global coordination.

This behavior corresponds directly to the algebra
of histories introduced earlier in the essay.
Execution extends histories through irreversible
events, replicas reconcile divergent branches
through join-like merge operations, and the
semilattice structure of the resulting histories
guarantees that the system converges toward a
consistent representation of the accumulated
execution trace.

\subsection{Version-Control Systems}

Distributed version-control systems provide a particularly
transparent illustration of the event-historical structure
described in the preceding sections. Systems such as Git
organize program development not around mutable states
but around an append-only history of commits that record
successive modifications to a codebase.

A repository in such a system is therefore most naturally
understood as a directed acyclic graph of commits, where
each commit represents a discrete event appended to the
history of the project. Rather than replacing previous
states, commits extend the historical trace of development,
preserving the causal relationships among modifications.
The present configuration of the repository is thus derived
from the accumulated sequence of historical events rather
than existing as an independent object detached from its
construction.

Branching within the repository corresponds to the
divergence of execution histories. Independent lines of
development extend a shared prefix of commits, producing
distinct historical trajectories that evolve concurrently.
Merge operations subsequently reconcile these trajectories
by constructing a new commit whose ancestry incorporates
both branches while preserving their causal ordering.

From the perspective developed in this essay, the commit
graph therefore realizes the same algebra of histories that
underlies the operational kernel. Commits correspond to
events extending the history, branches represent divergent
extensions of a shared prefix, and merges act as join-like
operations that reconcile compatible histories. The
structure of the repository is thus not simply a data
structure for storing code revisions but an explicit
manifestation of the event-historical semantics through
which the system evolves.

\subsection{Constraint-Based Programming}

Constraint-based programming languages provide another
illustration of the event-historical perspective developed
in this essay. Rather than specifying computation as a
sequence of explicit state transitions, such languages
describe systems in terms of relations that must hold
among variables. The role of execution is therefore not
to manipulate states directly but to progressively enforce
and reconcile these relations until a configuration
satisfying the specified constraints is obtained.

Operationally, execution proceeds through a sequence of
incremental updates in which local assignments are
introduced, revised, or propagated in response to the
constraints that govern the system. Each update modifies
the set of admissible configurations and therefore
reshapes the space of possible future transitions.
The system evolves through a process in which constraints
propagate through the network of variables until the
resulting configuration stabilizes.

Within the framework proposed here, this behavior can be
interpreted as a trajectory through the lattice of event
histories introduced earlier. Each constraint propagation
step corresponds to an event extending the history, while
the compatibility conditions imposed by the constraint
system determine which extensions are admissible. The
potential function associated with the constraint structure
guides execution toward configurations that minimize
inconsistency and satisfy the specified relations.

Seen from this perspective, constraint-based computation
does not simply search through a state space but instead
constructs a history of constraint-enforcing events whose
cumulative effects determine the eventual configuration
of the system. The final solution thus emerges as a stable
reduction of the underlying execution history rather than
as a primitive state reached through direct transitions.

\subsection{Physical Constraint Systems}

The same structural principles appear in physical
systems governed by local interactions.

In lattice models such as the Ising system,
each local update modifies the constraints
imposed on neighboring components.

Execution therefore proceeds through a sequence
of local events that reshape the space of
possible future transitions.

Stable configurations correspond to local minima
in the associated energy landscape.

\subsection{Common Structure}

Despite their apparent differences, these systems share
a common mathematical foundation. In each case the
underlying structure may be interpreted as a partially
ordered space of histories whose elements represent
irreversible sequences of events. Execution corresponds
to the monotonic extension of these histories, while
composition arises through merge operations that behave
as join constructions in a semilattice-like structure.
Abstraction, in turn, acts as a reduction mapping that
compresses histories into derived representations such
as states, summaries, or other forms of observational
projection.

Viewed in this way, the systems discussed above are not
isolated technical solutions but manifestations of the
same underlying algebraic organization. Histories grow
through extension, divergent branches are reconciled
through join-like merges, and the resulting structures
are interpreted through reductions that discard
information about the execution trace while preserving
those distinctions relevant to a given purpose.

This shared structure reveals a deep unity among systems
that are typically studied within separate domains of
computer science and physics. Distributed logs, version
control graphs, constraint-solving processes, and lattice
dynamics governed by local interactions all exhibit
variants of the same operational pattern: the monotonic
accumulation of events whose histories determine the
configuration of the system at any moment.

\subsection{Interpretation}

The event-historical kernel therefore provides
a minimal ontology of computation.

Rather than treating programs as transformations
between states, the kernel interprets computation
as the irreversible construction of histories
under constraint.

State, structure, and abstraction all emerge
as reductions of those histories.

This perspective aligns computational theory
with the dynamics of physical systems in which
structure arises from the cumulative effects
of locally constrained transitions.

\section{Execution as Ordered History Geometry}

The central claim of this paper is that computation is most naturally
understood not as the transformation of states but as the construction
and manipulation of histories.

A history records the irreversible accumulation of events. Each event
excludes alternative possibilities and therefore refines the set of
admissible continuations. Computation emerges from the lawful
composition of such exclusions.

Within this perspective, the minimal operational kernel consists of
three primitive operations: extension, merge, and reduction. Extension
adds new events to an existing history, producing longer executions.
Merge combines compatible histories by identifying shared prefixes and
reconciling their subsequent developments. Reduction removes distinctions
between histories by collapsing them into coarser summaries.

These operations generate an ordered structure over histories. The order
captures informational refinement: a history that records more events
contains strictly more constraints on the future. Composition of
histories corresponds to concatenation of exclusions, while merge and
reduction act as structural transformations within this ordered space.

The divisibility results discussed earlier illuminate a subtle property
of this geometry. When histories are represented operationally as
transformations acting on states, one obtains a particular notion of
lawful intermediate execution. When the same histories are represented
as transformations acting on observables, a dual notion emerges. The two
representations agree on empirical predictions but impose different
constraints on how histories may be factorized.

This reveals that execution possesses an intrinsic directionality. The
continuation of a history through extension is not equivalent to the
reconstruction of a history from its future boundary. In categorical
terms, histories admit two dual operational representations related by
opposition of composition order. The admissibility constraints governing
these representations need not coincide.

Consequently, the algebra of histories is not strictly self-dual. The
direction in which execution is evaluated—whether through forward
extension of states or backward propagation of observables—affects the
existence of lawful intermediate steps.

Seen from this perspective, computation resembles a geometric process
in an ordered space of histories. Extension traces directed paths
through this space, merge identifies compatible trajectories, and
reduction projects trajectories onto coarser descriptions. The
structure that emerges is closer to a geometry of execution than to a
traditional state-transition system.

Many familiar computational architectures implicitly operate within
this geometry. Distributed version control systems construct histories
through successive extensions and reconcile them through merge
operations. Event-sourced systems record append-only histories and
derive system state through reduction. Collaborative editing systems
maintain partially ordered logs of operations whose reconciliation
depends on compatibility conditions between histories.

The minimal operational kernel proposed here abstracts these patterns
into a general model of computation grounded in history composition.
Within this model, the existence of lawful intermediate histories
becomes a central structural property. The divisibility phenomena
observed in dynamical systems therefore appear not as isolated
physical curiosities but as manifestations of a deeper algebraic
feature of execution itself.

In this sense, computation may be understood as the geometry of
irreversible histories evolving within an ordered space of
possibilities.

\subsection{A Universal Property of the History Algebra}

The operations of extension, merge, and reduction generate the space of
executable histories in a precise sense. They form a minimal operational
kernel from which all history-based computations may be constructed.

Let $\mathcal{H}$ denote the algebra of histories generated by these
operations, subject to the compositional relations introduced earlier.
An execution in any concrete computational system may then be interpreted
as a structure-preserving map from $\mathcal{H}$ into the operational
algebra of that system.

More precisely, suppose $\mathcal{A}$ is a computational structure whose
executions admit operations corresponding to extension, merge, and
reduction. Then there exists a homomorphism

\[
\Phi : \mathcal{H} \to \mathcal{A}
\]

that preserves these operations.

\begin{proposition}[Universal property of the history algebra]
The history algebra $\mathcal{H}$ is initial among computational
structures equipped with extension, merge, and reduction operations.
That is, for any such structure $\mathcal{A}$ there exists a unique
structure-preserving map
\[
\Phi : \mathcal{H} \to \mathcal{A}.
\]
\end{proposition}

Intuitively, this property states that $\mathcal{H}$ contains no
structure beyond that required by the minimal operational kernel.
Any concrete model of computation implementing these operations
factors through the history algebra.

In this sense, the algebra of histories plays a role analogous to that
of a free algebra generated by a set of primitive operations. The
extension–merge–reduction kernel therefore provides a universal
description of history-based computation.

\section{Histories as an Order-Enriched Monoidal Category}

The structures introduced throughout this paper admit a concise
mathematical formulation.

Let $\mathcal{H}$ denote the category of executable histories. Objects
represent system interfaces or boundary conditions, while morphisms
represent history segments connecting those boundaries. Composition of
morphisms corresponds to temporal concatenation of histories.

\[
h_{u\to t} = h_{s\to t} \circ h_{u\to s}.
\]

In addition to composition, histories support a parallel combination
operation. Independent histories may be executed simultaneously, giving
rise to a monoidal product

\[
h_1 \otimes h_2.
\]

Together these operations endow $\mathcal{H}$ with the structure of a
monoidal category.

However, histories also carry an intrinsic informational ordering. A
history that records more events excludes more alternatives and
therefore contains strictly more information about the execution. This
naturally induces a partial order

\[
h_1 \preceq h_2
\]

whenever $h_2$ refines $h_1$ by recording additional events.

Composition and parallel combination are monotone with respect to this
ordering. If $h_1 \preceq h_2$ then

\[
f \circ h_1 \preceq f \circ h_2
\]

and

\[
h_1 \otimes g \preceq h_2 \otimes g
\]

whenever the compositions are defined.

Consequently the category of histories is naturally enriched over
ordered sets. Each hom-set $\mathcal{H}(x,y)$ forms a partially ordered
set describing the refinement structure of histories between those
boundaries.

Within this framework the primitive operations of the minimal
operational kernel acquire clear structural interpretations.

Extension corresponds to composition with additional history segments.
Merge corresponds to identifying compatible histories that share common
prefixes. Reduction corresponds to order-preserving maps that collapse
distinctions between histories.

Admissibility constraints restrict the morphisms that may appear in
valid executions. In many settings these constraints form a downward
closed subset of the hom-set order. The resulting structure describes
the space of lawful histories that a system may realize.

This categorical perspective clarifies why the direction of execution
matters. Representations of histories as state transformations and as
observable transformations correspond to functors

\[
F : \mathcal{H} \to \mathbf{OrdMon}, \qquad
G : \mathcal{H}^{op} \to \mathbf{OrdMon}.
\]

Although these representations are prediction-equivalent, they may
induce different admissibility constraints on intermediate history
segments. As a result, the factorization properties of histories can
depend on whether execution is evaluated through forward extension or
through its dual representation.

The algebra of histories therefore forms an order-enriched monoidal
structure whose operational interpretations need not be self-dual. This
structure provides the mathematical foundation for the history-based
model of computation developed in this paper.


% ------------------------------------------------------------
% 10. Conclusion
% ------------------------------------------------------------
\section{Conclusion}

This essay has proposed a minimal operational kernel for the
interpretation of computation in which the fundamental structure
of a system is not a static state but the irreversible history
of events through which the system has been constructed.
Traditional models of computation have generally described
programs as transformations between well-defined states.
While this perspective has yielded powerful formal tools,
it necessarily compresses the temporal and causal structure
of execution into a sequence of abstract snapshots.
Such compression obscures the processes by which systems
actually evolve and removes from view the irreversible
accumulation of events that produces observable behavior.

The framework developed here reverses this perspective by
treating event histories as the primary objects of computation.
Within this view, execution is understood as the monotonic
extension of histories through the occurrence of events,
composition arises through the reconciliation and merging
of compatible histories while preserving their causal
ordering relations, and abstraction appears as the reduction
of histories through the controlled loss of information.
These operations together form a minimal operational kernel
capable of generating the structures that appear across
contemporary computational systems.

When computation is interpreted in this manner, system states
cease to function as fundamental entities. Instead they emerge
as derived projections obtained by replaying or reducing the
underlying history of execution. Snapshots, summaries, and
other forms of representation therefore become informational
compressions of event traces rather than primary objects of
analysis. The apparent stability of a system state is thus
revealed to be a consequence of the structure of its history.

The algebra induced by these operations exhibits a natural
mathematical organization. Event histories form partially
ordered structures under prefix extension, admit join-like
merge operations when compatible branches of execution are
reconciled, and support reduction mappings that compress
histories into derived representations. This algebraic
structure is not unique to a particular programming paradigm.
Closely related forms appear in distributed log architectures,
version-control graphs, conflict-free replicated data types,
constraint-based programming systems, and even in physical
models such as lattice dynamics governed by local interaction
rules. In each of these domains the behavior of the system
emerges from the cumulative effects of locally constrained
events whose histories grow monotonically over time.

The significance of the event-historical kernel therefore
extends beyond a reformulation of programming semantics.
It suggests that computation, distributed coordination,
and certain classes of physical dynamical systems share a
common structural principle: complex behavior arises from
the irreversible accumulation of locally constrained
transitions whose histories determine the present
configuration of the system.

Under this interpretation, programs are no longer best
understood as static objects defined by membership within
sets of possible states. Instead they appear as processes
of construction unfolding through time, in which execution
records the history of events that produce structure,
composition reconciles divergent branches of that history,
and abstraction compresses the resulting traces into
representations suitable for reasoning and control.
State thus becomes a derived reduction of history,
structure emerges from the ordering relations among events,
and abstraction becomes the systematic compression of the
execution traces that give rise to the system itself.

\newpage
% ------------------------------------------------------------
% Bibliography
% ------------------------------------------------------------

\begin{thebibliography}{99}

\bibitem{lamport1978}
Leslie Lamport.
\textit{Time, Clocks, and the Ordering of Events in a Distributed System}.
Communications of the ACM, 21(7):558--565, 1978.

\bibitem{shapiro2011}
Marc Shapiro, Nuno Preguiça, Carlos Baquero, and Marek Zawirski.
\textit{Conflict-Free Replicated Data Types}.
In \textit{Stabilization, Safety, and Security of Distributed Systems (SSS)}, 2011.

\bibitem{kleppmann2017}
Martin Kleppmann.
\textit{Designing Data-Intensive Applications}.
O’Reilly Media, 2017.

\bibitem{fowler2005}
Martin Fowler.
\textit{Event Sourcing}.
martinfowler.com, 2005.

\bibitem{chacon2014}
Scott Chacon and Ben Straub.
\textit{Pro Git}.
Apress, 2014.

\bibitem{hopcroft2006}
John Hopcroft, Rajeev Motwani, and Jeffrey Ullman.
\textit{Introduction to Automata Theory, Languages, and Computation}.
Pearson, 2006.

\bibitem{winskel1993}
Glynn Winskel and Mogens Nielsen.
\textit{Models for Concurrency}.
In \textit{Handbook of Logic in Computer Science}. Oxford University Press, 1993.

\bibitem{winskel1987}
Glynn Winskel.
\textit{Event Structures}.
In \textit{Advances in Petri Nets 1986}, Lecture Notes in Computer Science 255:325--392.
Springer, 1987.

\bibitem{settimo2026}
Federico Settimo, Andrea Smirne, Kimmo Luoma, Bassano Vacchini,
Jyrki Piilo, and Dariusz Chru\'sci\'nski.
\textit{Divisibility of Dynamical Maps: Schr{\"o}dinger Versus Heisenberg Picture}.
PRX Quantum, 7(1), 2026.
\url{https://doi.org/10.1103/6dt2-sq44}

\bibitem{breuer2009}
Heinz-Peter Breuer, Elsi-Maria Laine, and Jyrki Piilo.
\textit{Measure for the Degree of Non-Markovian Behavior of Quantum Processes in Open Systems}.
Physical Review Letters, 103(21):210401, 2009.

\bibitem{rivas2010}
Ángel Rivas, Susana F. Huelga, and Martin B. Plenio.
\textit{Entanglement and Non-Markovianity of Quantum Evolutions}.
Physical Review Letters, 105(5):050403, 2010.

\bibitem{sipser2013}
Michael Sipser.
\textit{Introduction to the Theory of Computation}.
Cengage Learning, 2013.

\bibitem{ising1925}
Ernst Ising.
\textit{Contribution to the Theory of Ferromagnetism}.
Zeitschrift für Physik, 31:253--258, 1925.

\bibitem{brush1967}
Stephen G. Brush.
\textit{History of the Lenz–Ising Model}.
Reviews of Modern Physics, 39(4):883--893, 1967.

\bibitem{mezard1987}
Marc Mézard, Giorgio Parisi, and Miguel Virasoro.
\textit{Spin Glass Theory and Beyond}.
World Scientific, 1987.

\bibitem{mackay2003}
David J. C. MacKay.
\textit{Information Theory, Inference, and Learning Algorithms}.
Cambridge University Press, 2003.

\end{thebibliography}

\end{document} 